\clearpage\fancyhead[L]{\small PALLAS.}\hypertarget{native-354-lower}{}\label{native-354-lower}
\begin{center}\textbf{PALLAS.}\par\vspace{4pt}\hrule\vspace{4pt}\fontsize{10}{12}\selectfont Göttingische gelehrte Anzeigen. Stück 67. Seite 657\;..\;660. 1812 April 25.\par\vspace{4pt}\hrule\vspace{4pt}\end{center}

\par\indent Schon einige Male ist in unsern Blättern von der Arbeit des Prof. \textsc{Gauss} über die Störungen der Pallas die Rede gewesen (man s. die Stücke 8 und 130 vom vor. J. [S. 322 und 327 d. B.]). Die dort angedeuteten Resultate dieser Untersuchungen haben jetzt abermals eine schöne Bestätigung erhalten. Am 4. April erlaubte das Wetter zum ersten Mal in diesem Jahre, den Planeten aufzusuchen. Genau auf dem Platze, den die vom Hrn. \textsc{Nicolai} nach der Störungstheorie berechnete Ephemeride anwies, war ein Sternchen eilfter Grösse sichtbar, welches am 9. April diesen Platz verlassen, und sich, wie gehörig, fortbewegt hatte. Ein paar Vergleichungen mit 93 Herculis gaben folgende Stellung, wobei indessen die Declination blos geschätzt ist.

\begin{center}\fontsize{10}{12}\selectfont\setlength{\tabcolsep}{0pt}\renewcommand{\arraystretch}{1.15}\begin{tabular}{l}
April 9. $11^{\mathrm u}\;32^{\mathrm m}\;27^{\mathrm s}$ M. Z.\\
Gerade Aufsteigung der Pallas $269^\circ\;6'\;6''7$\\
Nordliche Abweichung $16\;38$
\end{tabular}\end{center}
\SourceLeaf{355}\Needspace{13\baselineskip}\par\indent Hienach stimmt der Stand des Planeten vollkommen mit jener Ephemeride überein. Gleich nachher erhielt Prof. \textsc{Gauss} vom Hrn Dr. \textsc{Olbers}, welchen er zur Aufsuchung des Planeten aufgefordert hatte, die Nachricht, dass ihm diese schon am 3. April geglückt sei. Die beiden folgenden Bremer Beobachtungen stimmen gleichfalls aufs genaueste mit der Ephemeride des Hrn. \textsc{Nicolai}:

\begin{center}
\fontsize{10}{12}\selectfont\setlength{\tabcolsep}{3pt}\renewcommand{\arraystretch}{1.15}\begin{tabular}{c|c|c}
Mittlere Zeit & \shortstack{Gerade\\Aufsteigung} & \shortstack{Nordliche\\Abweichung} \\
\hline
April 3 \quad $11^{\mathrm u}\;26^{\mathrm m}\;44^{\mathrm s}$ & $268^\circ\;42'\;23''1$ & $15^\circ\;25'\;4''9$ \\
\phantom{April} 4 \qquad\quad $12\;2\;33$ & $268\;47\;12.5$ & $15\;37\;41.1$ \\
\end{tabular}
\end{center}

\par\indent Wir eilen, diese Stellungen des durch seine Lichtschwäche in diesem Jahre äusserst schwer zu beobachtenden Planeten bekannt zu machen, und fordern alle Astronomen angelegentlich auf, ihn besonders um die Zeit der Opposition (10. Juni) sorgfältig zu observiren.

\par\indent Auch die \emph{erste vorläufige} Berechnung der \emph{allgemeinen} Theorie der Pallasstörungen hat Prof. \textsc{Gauss} seit kurzem vollendet. Die Anzahl aller periodischen Gleichungen von der Einwirkung Jupiters, die noch merklich sind, steigt an \emph{Vier hundert}. Unter den sehr merkwürdigen, schon hieraus fliessenden, Resultaten ist besonders eins von höchstem Interesse. Aus Gründen legen wir es hier in folgender Chiffre nieder, wozu wir zu seiner Zeit den Schlüssel geben werden:

\begin{center}
1111000100101001
\end{center}

\par\indent Wir fügen bei dieser Gelegenheit noch zwei andere interessante astronomische Mittheilungen bei, aus einem Briefe des Hrn. Prof. \textsc{Bessel} in Königsberg an den Prof. \textsc{Gauss}, vom 26. März.

\par\indent Bekanntlich hat schon \textsc{Herschel} die Meinung aufgestellt, dass diejenigen Sterne, welche uns als Doppelsterne erscheinen, nicht durch Zufall mit uns in Einer geraden Linie, sondern wirklich nahe beisammen stehen, und jedes Paar ein eigenes System ausmachen. Wenn auch \emph{diejenigen Gründe}, aus welchen dieser berühmte Beobachter diesen Schluss zog, noch erheblichen Einwürfen unterliegen, so sprechen doch andere Beobachtungen sehr stark für die Meinung selbst. Es ist nemlich aus Gründen der Wahrscheinlichkeitsrechnung höchst unwahrscheinlich, dass \emph{der Zufall} aus den Sternen bis zu einer gewissen Ordnung auch nur ein einziges Paar Doppelsterne gebildet haben sollte, geschweige denn so viele, als sich wirklich am Himmel vorfinden. Jetzt hat aber diese Meinung durch eine sehr interessante Bemerkung des Hrn. Prof. \textsc{Bessel} eine sehr einleuchtende Bestätigung erhalten. Seine wichtige Bearbeitung der \textsc{Bradley}ischen Beobachtungen, mit welcher er sich schon seit mehreren Jahren beschäftigt, hat ihm mehrere Fixsterne kenntlich gemacht, die eine sehr starke eigene Bewegung haben. Unter diesen ist am merkwürdigsten der Doppelstern 61 Cygni, dessen eigene Bewegung von allen bekannten die stärkste ist. Diese Bewegung ist beiden Sternen gemeinschaftlich, und ihre Richtung auf den Hercules zu, welcher Richtung auch alle starken eigenen Bewegungen sehr nahe folgen. Für 1755 findet Hr. \textsc{Bessel} aus \textsc{Bradley}'s Beobachtungen:

\begin{center}
\fontsize{10}{12}\selectfont\setlength{\tabcolsep}{3pt}\renewcommand{\arraystretch}{1.15}\begin{tabular}{c|c|c}
 & \shortstack{Gerade\\Aufsteigung} & \shortstack{Nordliche\\Abweichung} \\
\hline
61 Cygni & $313^\circ\;59'\;11''3$ & $37^\circ\;33'\;29''7$ \\
Sequens & $313\;59\;25.7$ & $47\;33\;45.7$ \\
\end{tabular}
\end{center}

\par\indent \textsc{Piazzi}'s Angabe konnte nicht benutzt werden, weil das Jahr der Beobachtung unbekannt ist; allein zwei Beobachtungen der \emph{Histoire céleste} geben die jährliche Bewegung des Hauptsterns

\begin{center}\fontsize{10}{12}\selectfont\setlength{\tabcolsep}{0pt}\renewcommand{\arraystretch}{1.15}\begin{tabular}{l}
$+5''250$ in gerader Aufsteigung;\\$+2''321$ in Declination.
\end{tabular}\end{center}
\SourceLeaf{356}\noindent Könnte man die jährliche Parallaxe dieses Doppelsterns durch Beobachtungen ausmitteln, so würde in Zukunft, sobald die Bewegung der beiden Sterne um ihren gemeinschaftlichen Schwerpunkt bekannt sein wird, das Verhältniss der Summe ihrer Massen zu der Masse unserer Sonne daraus abgeleitet werden können. Man sieht hieraus, wie sehr dieser merkwürdige Stern die höchste Aufmerksamkeit der Astronomen verdient.

\Needspace{25\baselineskip}\par\indent Eine zweite Mittheilung des Hrn. Prof. \textsc{Bessel} betrifft den grossen Cometen des vorigen Jahres. Schon vor acht Monaten (s. unsere gel. Anz. vom vor. J. S. 1293 [S. 329 d. B.]) bemerkten wir, dass Hoffnung da sei, diesen Cometen noch zum dritten Male im April und den folgenden Monaten im Wassermann wieder auffinden zu können. Man hat noch keine Ursache, diese Hoffnung aufzugeben, da der Comet unter der Voraussetzung, dass seine Helligkeit nicht durch physische Ursachen modificirt werde, doch noch ein Fünftel so viel Licht haben wird, als zu Anfang Januars, wo man ihn noch mit blossen Augen erkannte. Folgende kleine, vom Hrn Prof. \textsc{Bessel} berechnete, Ephemeride wird bei dieser Aufsuchung gute Dienste leisten:

\begin{center}
\fontsize{10}{12}\selectfont\setlength{\tabcolsep}{3pt}\renewcommand{\arraystretch}{1.15}\begin{tabular}{c|c|c|c}
\shortstack{M. Pariser Zeit\\$1812\;16^{\mathrm u}$} & \shortstack{Gerade\\Aufsteig.} & \shortstack{Südliche\\Abweich.} & Lichtstärke \\
\hline
April 12 & $335^\circ\;44'$ & $6^\circ\;52'$ & 0.0068 \\
16 & $336\;14$ & $7\;6$ & 0.0067 \\
20 & $336\;42$ & $7\;21$ & 0.0065 \\
24 & $337\;6$ & $7\;37$ & 0.0064 \\
28 & $337\;28$ & $7\;55$ & 0.0063 \\
Mai\ \ 2 & $337\;46$ & $8\;14$ & 0.0062 \\
6 & $338\;1$ & $8\;35$ & 0.0061 \\
10 & $338\;12$ & $8\;57$ & 0.0061 \\
14 & $338\;20$ & $9\;21\tfrac{1}{2}$ & 0.0060 \\
18 & $338\;23$ & $9\;48$ & 0.0060 \\
22 & $338\;23$ & $10\;16$ & 0.0059 \\
26 & $338\;18$ & $10\;46$ & 0.0059 \\
30 & $338\;10$ & $11\;19$ & 0.0058 \\
\end{tabular}
\end{center}
